\booksection{第十一章附：專業影像技術知識}

\chapterintro{想要真正理解手機攝影和影片的奧秘，你需要了解一些專業的影像技術知識。這些知識不僅能幫你更好地使用手機相機，還能讓你在選擇裝置時更有判斷力。}

\booksubsection{一、解析度、幀率與位元速率：影片質量的三要素}

\subsubsection*{解析度（Resolution）：}

解析度決定了畫面的清晰度，用畫素數表示。

\begin{itemize}
  \item \textbf{720P (HD)}：1280×720畫素，入門級清晰度。
  \item \textbf{1080P (Full HD)}：1920×1080畫素，目前最主流的解析度。
  \item \textbf{2K (QHD)}：2560×1440畫素，比1080P更清晰。
  \item \textbf{4K (UHD)}：3840×2160畫素，超高畫質，細節豐富。
  \item \textbf{8K}：7680×4320畫素，未來標準，目前裝置支援有限。
\end{itemize}

解析度越高，畫面越清晰，但檔案也越大，對儲存和處理能力要求越高。

\subsubsection*{幀率（Frame Rate）：}

幀率決定了畫面的流暢度，用每秒幀數（fps）表示。

\begin{itemize}
  \item \textbf{24fps}：電影標準幀率，有獨特的"電影感"。
  \item \textbf{30fps}：電視和網路影片常用，流暢自然。
  \item \textbf{60fps}：高幀率，適合運動場景和遊戲錄製，非常流暢。
  \item \textbf{120fps/240fps}：超高幀率，可用於慢動作影片。
\end{itemize}

幀率越高，畫面越流暢，慢動作效果越好，但檔案也越大。

\subsubsection*{位元速率（Bitrate）：}

位元速率決定了影片的資料密度，用每秒資料量（Mbps）表示。

\begin{itemize}
  \item 位元速率越高，畫質越好，細節保留越多，但檔案也越大。
  \item 位元速率過低會導致畫面出現馬賽克、色塊等壓縮痕跡。
  \item 專業影片拍攝通常需要高位元速率以保留更多後期處理空間。
\end{itemize}

\booksubsection{二、影片編碼格式：H.264、H.265與ProRes}

\subsubsection*{H.264 (AVC)：}

\begin{itemize}
  \item \textbf{特點}：最廣泛使用的編碼格式，相容性極好。
  \item \textbf{優點}：幾乎所有裝置都能播放，壓縮效率適中。
  \item \textbf{缺點}：相比新格式，壓縮效率較低，檔案較大。
  \item \textbf{適用場景}：網路影片分享、普通使用者日常拍攝。
\end{itemize}

\subsubsection*{H.265 (HEVC)：}

\begin{itemize}
  \item \textbf{特點}：H.264的升級版，壓縮效率提升約50\%。
  \item \textbf{優點}：相同畫質下檔案更小，支援更高解析度和幀率。
  \item \textbf{缺點}：編碼和解碼需要更強的處理器，部分老裝置不相容。
  \item \textbf{適用場景}：4K影片拍攝、儲存空間有限的場景。
\end{itemize}

\subsubsection*{Apple ProRes：}

\begin{itemize}
  \item \textbf{特點}：蘋果開發的專業影片編碼，幾乎無失真壓縮。
  \item \textbf{優點}：畫質極高，後期處理空間大，色彩準確。
  \item \textbf{缺點}：檔案極大，需要大量儲存空間，處理需要專業裝置。
  \item \textbf{適用場景}：專業影片製作、電影級拍攝。
\end{itemize}

\subsubsection*{ProRes RAW：}

\begin{itemize}
  \item \textbf{特點}：結合了ProRes和RAW格式的優勢。
  \item \textbf{RAW是什麼}：未經處理的原始感測器資料，保留最大資訊量。
  \item \textbf{優點}：後期調整空間極大，可以改變曝光、白平衡等引數而不損失畫質。
  \item \textbf{缺點}：檔案極其巨大，需要專業後期軟體處理。
  \item \textbf{適用場景}：高階專業製作。iPhone 15 Pro系列及以後支援拍攝Apple ProRes影片（最高4K 60fps），iPhone 17 Pro進一步支援ProRes RAW直接機內錄製。
\end{itemize}

\booksubsection{三、HDR：高動態範圍技術}

HDR（High Dynamic Range）技術可以同時保留畫面中最亮和最暗部分的細節。

\begin{itemize}
  \item \textbf{HDR10}：最基礎的HDR標準，相容性好。
  \item \textbf{HDR10+}：三星推動的標準，支援動態後設資料。
  \item \textbf{Dolby Vision}：杜比開發的頂級HDR標準，畫質最佳，但需要授權。
  \item \textbf{HLG}：廣播級HDR標準，相容普通顯示器。
\end{itemize}

拍攝HDR影片時，需要支援HDR的顯示器才能看到完整效果。

\booksubsection{四、色彩空間與色深}

\subsubsection*{色彩空間：}

\begin{itemize}
  \item \textbf{sRGB}：標準色彩空間，相容性最好。
  \item \textbf{Display P3}：蘋果推廣的廣色域標準，色彩更豐富。
  \item \textbf{Rec.2020}：HDR影片的標準色彩空間，色域最廣。
\end{itemize}

\subsubsection*{色深：}

\begin{itemize}
  \item \textbf{8-bit}：1670萬色，普通影片標準。
  \item \textbf{10-bit}：10.7億色，專業影片標準，色彩過渡更平滑。
  \item \textbf{12-bit}：680億色，高階專業標準。
\end{itemize}

\begin{tcbinfo}{實用建議}
\begin{itemize}
  \item 日常拍攝：1080P/4K + 30fps + H.265編碼，平衡畫質和檔案大小。
  \item 運動拍攝：1080P/4K + 60fps，保證流暢度。
  \item 專業創作：4K + ProRes（如果支援），保留最大後期空間。
  \item 儲存建議：高位元速率影片很佔空間，建議定期備份到電腦或雲盤。
\end{itemize}

\end{tcbinfo}

\begin{tcbpractice}{去探索一下：你的手機影片設定}
\begin{enumerate}
  \item 開啟相機App，切換到影片模式。
  \item 檢視解析度和幀率選項，嘗試不同的組合。
  \item 如果支援，檢視是否有HDR、ProRes等高階選項。
  \item 拍攝幾段不同設定的影片，對比畫質和檔案大小。
\end{enumerate}

\end{tcbpractice}

\booksubsection{五、RAW格式與HEIF/JPEG對比}

\subsubsection*{什麼是RAW格式：}

RAW格式是相機感測器記錄的“原始資料”——未經任何壓縮和處理的純原始資訊。它就像生鮮食材，需要後期“烹飪”才能成品。手機因為儲存和算力有限，預設拍出來的都是已經“烹飪好”的JPEG或HEIF，RAW需要專門開啟。

\subsubsection*{RAW vs HEIF vs JPEG對比：}

\begin{itemize}
  \item \textbf{JPEG}：最老的格式，相容性最好（任何裝置都能開啟）。8-bit色深，有失真壓縮，檔案最小（一張照片2-5MB）。後期空間有限，強行提亮暗部容易出現色塊。
  \item \textbf{HEIF/HEIC}：蘋果主推的新格式，同樣畫質下檔案只有JPEG的一半（一張照片1-3MB）。支援10-bit色深和透明通道。缺點是部分老裝置和Windows需要安裝外掛才能開啟。
  \item \textbf{RAW}：未經壓縮的原始資料，檔案最大（一張照片20-100MB）。色深通常12-16bit，後期空間最大。缺點是必須用專業軟體（Lightroom、Snapseed等）才能開啟和處理，無法直接分享。
\end{itemize}

\subsubsection*{為什麼專業攝影師要拍RAW：}

\begin{itemize}
  \item \textbf{白平衡可後期無損調整}：JPEG拍錯了白平衡（比如白熾燈下偏黃），後期糾正會損失畫質；RAW可以重新設定色溫，完全無損。
  \item \textbf{曝光容差大}：JPEG欠曝2檔的暗部基本是死黑一片，提亮後全是噪點；RAW能從欠曝3-4檔的暗部"救"出可用細節。過曝同理。
  \item \textbf{暗部和高光細節更豐富}：RAW的動態範圍遠大於JPEG，亮部不過曝、暗部不死黑的餘地大得多。
  \item \textbf{銳化、降噪、鏡頭矯正都在後期可控}：相機的機內處理是"一錘子買賣"，RAW可以反覆調整，每次都從原始資料出發。
\end{itemize}

\subsubsection*{手機如何開啟RAW：}

\begin{itemize}
  \item \textbf{iPhone ProRAW}：僅iPhone 12 Pro及以上Pro機型支援。開啟"設定 相機 格式"中開啟"Apple ProRAW"。然後在相機介面頂部會出現一個"RAW"按鈕，點亮即可。檔案約25MB/張，建議有需求時再開。
  \item \textbf{安卓Motion Raw / Ultra RAW}：各品牌叫法不同。小米/OPPO/vivo旗艦機型在相機"專業模式"或"更多"裡能找到RAW開關；華為在"專業模式"中開啟RAW；三星Galaxy S系列在"專業"模式裡有RAW選項。檔案大小通常在20-50MB。
  \item \textbf{通用建議}：日常隨手拍沒必要開RAW，檔案太大佔空間；只有在重要場景（風光、人像、需要後期）時才開啟。
\end{itemize}

\booksubsection{六、計算攝影原理}

現代手機攝影已經離不開“計算攝影”——也就是靠演算法而不僅是鏡頭和感測器來生成照片。理解這幾個核心原理，能讓你更明白手機在按下快門的瞬間做了什麼。

\subsubsection*{多幀合成：}

手機按下一次快門，其實拍了多張照片。演算法把多幀對齊、加權融合，得到一張高質量成片。

\begin{itemize}
  \item \textbf{夜景模式}：1-3秒內連拍多張欠曝照片，對齊後疊加——噪點互相抵消、亮度累加。結果就是暗光也能拍出乾淨的畫面。
  \item \textbf{HDR（高動態範圍）}：分別拍攝欠曝、正常、過曝三張（或多張），合成一張亮部和暗部都有細節的照片。逆光場景尤其有效。
  \item \textbf{Deep Fusion（蘋果）/Photonic Engine（光子引擎）}：針對中低光場景，逐畫素分析多幀並融合，重點提升紋理和細節，毛衣、毛髮這類細節會明顯更清晰。
\end{itemize}

\subsubsection*{語義分割：}

AI不只是把整張圖一刀切處理，而是先“看懂”畫面裡有什麼——天空、人臉、草地、建築、衣服——然後對每個區域分別應用不同的處理策略。

\begin{itemize}
  \item \textbf{天空}：單獨增強藍色、壓低亮度，避免過曝。
  \item \textbf{人臉}：識別膚色區域，單獨提亮、平滑膚色，不動五官和頭髮細節。
  \item \textbf{草地/樹葉}：增強綠色飽和度，提升細節銳化。
  \item \textbf{背景虛化}：人像模式下，AI識別主體輪廓，背景模糊、主體清晰。
\end{itemize}

\subsubsection*{深度學習的應用：}

\begin{itemize}
  \item \textbf{AI場景識別}：相機自動識別"美食""寵物""夜景""文件"等場景，並對應最佳化引數。比如美食模式會增強暖色調和飽和度，讓食物看起來更誘人。
  \item \textbf{AI人像虛化}：手機鏡頭物理上無法做到單反那種淺景深，靠深度學習模型預測主體邊緣，再對背景做模糊。最新演算法已經能正確處理頭髮絲、玻璃邊緣這些細節。
  \item \textbf{AI消除/AI擴圖}：選中畫面中的路人或雜物，AI基於上下文"補全"被遮擋的部分。這是生成式AI在影像領域的應用。
\end{itemize}

\subsubsection*{計算攝影 vs 光學攝影：}

\begin{itemize}
  \item \textbf{計算攝影的優勢}：小體積也能得到好畫質，暗光、逆光等極端場景也能出片，普通使用者不需要懂攝影知識就能拍出"好看"的照片。
  \item \textbf{計算攝影的劣勢}：畫面有時"過度處理"——銳化過度、色彩失真、天空顏色不自然；多幀合成在主體移動時會留下重影；細節是"算出來的"而非真實記錄的。
  \item \textbf{光學攝影的優勢}：大底感測器+大光圈鏡頭能真實捕捉光線，畫質自然，後期空間大；不會因為演算法"腦補"產生偽影。
  \item \textbf{趨勢}：現代旗艦手機是"光學+計算"結合——用大底感測器獲取真實資訊，用演算法最佳化呈現。兩者並不對立。
\end{itemize}

\booksubsection{七、鏡頭型別與焦段}

現代旗艦手機通常配三到四顆鏡頭，每顆對應不同焦段。理解焦段，能讓你在不同場景下選對鏡頭。

\subsubsection*{常見鏡頭與等效焦距：}

\begin{itemize}
  \item \textbf{超廣角（13-24mm等效）}：視野最廣，適合風光、建築、室內空間。能拍到眼前所見的大場面，但邊緣會有畸變（直線變彎），拍人像時人臉會被拉伸。適合"大場景"，不適合"特寫"。
  \item \textbf{主攝（24-28mm等效）}：日常最常用的鏡頭，畫質最好、光圈最大。兼顧視野和畫質，拍人拍景都能用。大多數手機的"1x"就是主攝。
  \item \textbf{長焦（70-200mm等效）}：用於拍遠景、人像特寫。長焦能壓縮空間感（遠處的山和近處的人看起來更近），拍人像能獲得更自然的五官比例和背景虛化。
  \item \textbf{潛望式長焦（200mm+等效）}：透過反射式光路設計，在手機厚度內實現更長焦距。能在很遠距離拍清楚細節，適合演唱會、體育賽事、野生動物。旗艦機型才有。
\end{itemize}

\subsubsection*{感測器尺寸與畫素大小的權衡：}

感測器是“底”，決定能接收多少光線。行業裡說“底大一級壓死人”——同樣畫素數下，感測器越大，單個畫素也越大，進光量越多，畫質越好。

\begin{itemize}
  \item \textbf{大底感測器}：進光量大，夜景畫質好，動態範圍寬。缺點是手機要做厚（鏡頭模組凸起），成本高。
  \item \textbf{高畫素小底}：白天光線充足時，高畫素能拍出更清晰的照片，方便後期裁切。但暗光下每個畫素進光不足，噪點多。
  \item \textbf{主流旗艦選擇}：主攝多為1/1.3英寸到1英寸（部分機型已用上1英寸大底），畫素數5000萬左右——這是畫質和體積的平衡點。
\end{itemize}

\subsubsection*{畫素融合技術（Pixel Binning）：}

為了兼顧白天高畫素和夜晚大畫素，手機廠商用“畫素融合”——把多個小畫素合併成一個大畫素輸出。

\begin{itemize}
  \item \textbf{四合一畫素（Quad Bayer）}：5000萬畫素感測器，輸出1250萬畫素照片。每個輸出畫素由4個相鄰畫素合成，進光量提升4倍，畫質更乾淨。
  \item \textbf{九合一畫素（Nonacell）}：用於1億畫素以上的感測器，輸出約1100-1200萬畫素。同樣原理，進光量更大。
  \item \textbf{實際效果}：預設拍照模式下，你拿到的就是融合後的低畫素照片——畫質好、檔案小。需要極致細節時（比如拍文件、遠景），可以手動切到"高畫素模式"輸出全解析度。
\end{itemize}

\booksubsection{八、特殊拍攝模式}

除了日常拍照，手機裡還有幾個特殊模式，用對了能拍到平時拍不出的畫面。

\subsubsection*{全景模式：}

\begin{itemize}
  \item 拍攝時按提示箭頭緩慢平穩地橫向移動手機，不要快——快了會出現拼接錯位。
  \item 保持手機垂直方向穩定，只橫向平移。如果上下晃動，地平線會變波浪。
  \item 有移動的人或車時慎用——拼接時會出現"半個身體"或"鬼影"。
  \item 全景模式也可以豎著拍——豎拍橫拼接，適合拍高聳的建築或瀑布。
\end{itemize}

\subsubsection*{慢動作與延時攝影：}

\begin{itemize}
  \item \textbf{慢動作}：用120fps或240fps高幀率拍攝，回放時變成慢鏡頭。適合拍水滴、運動瞬間、寵物跳躍。光線要充足，慢動作模式下進光時間短，暗光畫質差。
  \item \textbf{延時攝影}：間隔拍攝多張照片合成快進影片。適合拍車流、雲彩移動、日落、花朵開放。拍攝時長幾分鐘到幾小時不等，務必用三腳架固定手機，手持必糊。
  \item \textbf{小技巧}：延時攝影拍雲彩時，把手機架好，拍15-30分鐘，最後得到10秒左右的延時影片，效果震撼。
\end{itemize}

\subsubsection*{星空/天文模式：}

\begin{itemize}
  \item 需要遠離城市光汙染，找完全沒有路燈干擾的郊外或山區。城市裡拍不出銀河。
  \item 必須用三腳架，手持必糊。手機要穩穩架住幾秒到幾十秒。
  \item 天氣要晴朗無雲，月亮不要太大（滿月會蓋住星光）。農曆月初或月末的暗夜最適合。
  \item iPhone的"夜間模式"會自動識別星空並延長曝光時間（最長30秒）；安卓旗艦的"星空模式"或"天文模式"效果類似。要耐心等待，按下快門後不要動手機。
\end{itemize}

\subsubsection*{微距模式：}

\begin{itemize}
  \item 用於拍攝花朵、昆蟲、珠寶、文字特寫等小物體。能貼近到1-2釐米距離對焦。
  \item 部分手機有專門的微距鏡頭（多為200萬畫素或500萬畫素的小底鏡頭），畫質一般；旗艦機用超廣角鏡頭兼顧微距，畫質更好。
  \item 貼近拍攝時鏡頭離主體很近，會遮擋光線，必要時用檯燈或手電補光。
  \item 對焦距離極近，景深很淺，主體前後移動幾毫米就可能脫焦。拍攝時手要穩，最好讓主體保持不動。
\end{itemize}

\begin{tcbinfo}{靜態攝影實用建議}
\begin{itemize}
  \item 日常拍攝：HEIF/JPEG預設即可，檔案小、畫質夠用。
  \item 重要場景（風光、人像、需要後期）：開啟RAW，保留最大後期空間。
  \item 選對鏡頭：風光用超廣角，人像用長焦，日常用主攝。
  \item 夜景和星空：務必固定手機，開夜景模式或天文模式。
  \item 大底手機（如1英寸感測器機型）畫質上限更高，但操作習慣仍是關鍵。
\end{itemize}

\end{tcbinfo}

\begin{tcbpractice}{動手試試：摸清你的相機能力}
\begin{enumerate}
  \item 檢視你的手機是否支援RAW拍攝（設定相機格式），試著開啟並拍攝一張，對比同一場景的JPEG和RAW檔案大小。
  \item 切換到超廣角、主攝、長焦三個焦段，分別拍同一處遠景，體會視角差異和畫質差異。
  \item 找一個有夜景模式或星空模式的場景，固定手機後拍攝一張夜景照片。
  \item 嘗試用全景模式橫拍一張大場景，注意保持手機平穩移動。
\end{enumerate}

\end{tcbpractice}
